\section{技术补充与复现细节}
\label{app:technical}
\label{app:implementation}

本附录说明正文模型如何转换为可重复执行的计算过程。公式中的计划购电量$G_t^0,G_t$、充放电量$C_t,D_t$、紧急购电量$X_t$及余电量$W_t$均为10分钟区间的交流侧电量，储电量记为$E_t$；单向效率为$\eta=0.9$。完整程序紧接本附录列出，运行顺序与依赖版本见支撑材料\texttt{README.md}。

\subsection{问题二风险裕度—即时反馈模型的实现}
问题二正式入口为\texttt{simulate\_quantile\_q2}。每天只拟合一次预测权重：前一日、前一周和最近7日中位数构成三列基预测，最近14日内全部10分钟样本共同最小化净负荷绝对误差；负载权重和光伏权重分别非负且和为1。历史不足8日时采用前一日曲线，首日采用附件1。最近30日已揭晓净负荷残差用于计算70\%、80\%或90\%分位数；逐段分位数与同小时池化分位数各取一半。次日预测以当日预测代替未揭晓的当日实际曲线，故48小时输入不含未来实测值。

48小时稀疏线性规划以计划购电费为目标，固定起点真实储电量和末端6000 kWh，仅提交前144段。日内每10分钟直接按当前实际余缺执行：余量为正时充到功率或容量边界，余量为负时先放到功率或储电下界，仍有缺口才记紧急购电；不调用另一个MPC函数。配置字段\texttt{model=rolling\_quantile}、\texttt{feedback=immediate\_safe}和\texttt{control\_hours=0}共同标识该正式路径。

70\%、90\%、无分位数裕度和无储能方案均使用同一实现，只改变一个预先声明的参数。作为并列候选的12情境随机规划方案由\texttt{scenario\_reference}独立运行，其预测审计与逐日结果置于\texttt{q2\_model\_comparisons/}；完整逐时账本可由同一入口重建。

\subsection{12情境模型及问题三、四的预测与情景生成}
\label{app:forecast}

\paragraph{26维特征构造。}
为说明特征的时间关系，令$D$为本次午夜预测的日序号，$a$为目标日序号，$s=0,\ldots,143$为日内区间。把三个预测目标写为
\begin{equation}
 z_{a,s}=(L_{a,s},V_{a,s},\log p_{a,s})^{\mathsf T},\qquad
 A=\operatorname{diag}(5000,5000,1).
\end{equation}
其中$L,V$单位为kW，$A$用于统一残差的数值尺度。朴素预测的负载、光伏基线见正文；对数价格基线依次以$0.60,0.30,0.10$加权过去28天同星期均值、最近7天均值和附件1先验。历史不足时取已有天数；无同星期样本时改用最近样本；首日直接使用附件1。

每个目标区间的输入$x_{D,a,s}$含四组各3维的曲线特征：基线、起点前最近完整日、目标日的7天滞后、附件1先验，均左乘$A^{-1}$。最近完整日固定为$D-1$；只有$0\le a-7<D$时才取实际7天滞后，否则使用先验。其余14维为
\begin{equation}
 (\sin\theta_s,\cos\theta_s,\sin2\theta_s,\cos2\theta_s,
   \sin\phi_a,\cos\phi_a,a-D,\text{目标星期的7维独热编码}),
\end{equation}
其中$\theta_s=2\pi(s+1)/144$，$\phi_a=2\pi a/365.25$。
独热编码是以一个1和六个0表示星期类别。预测次日时，最近日特征仍取$D-1$，不会改取尚未结束的$D$日。固定电价问题在全部特征和残差计算中使用附件1日价格曲线；附件4实际价格仅进入问题四。

\paragraph{岭回归与梯度提升残差校正。}
设训练目标为$y_i=A^{-1}(z_i-\widehat z_i^{\rm base})$，以$y_i^{\mathsf T}$为行组成矩阵$Y$，列向量$\mu_Y$记录三个目标的训练均值。岭回归先用训练样本均值和标准差标准化$X$，标准差小于$10^{-12}$的列以1替代除数，再求解
\begin{align}
 \widehat W&=\arg\min_W\left\|Y-\bm1\mu_Y^{\mathsf T}-X_{\rm std}W\right\|_F^2
                       +30\|W\|_F^2,\\
 \widehat z&=\widehat z^{\rm base}+A(\mu_Y+\widehat W^{\mathsf T}x_{\rm std}).
\end{align}
截距不受惩罚。实现求解加上$30I$的正规方程，并检查系数有限性。

提升树分别拟合三个残差目标。以某一目标为例，从训练均值$F_0$开始，用回归树逐轮拟合尚未解释的残差：
\begin{equation}
 F_m(x)=F_{m-1}(x)+0.08\,g_m(x),\qquad m=1,\ldots,60.
\end{equation}
对于给定的树分区，记所拟合目标的残差为$y_i^{\ast}$，叶值对应平方损失与L2惩罚
$\frac12\sum_i\{y_i^{\ast}-F_{m-1}(x_i)-g_m(x_i)\}^2+\frac52\sum_j b_j^2$，
其中$b_j$为叶节点的输出；分裂由直方图算法搜索。固定最多15个叶节点、每叶最少36个样本、L2参数5，随机种子2026，关闭早停。树的输出同样乘回$A$并加到基线上。

\paragraph{滚动训练与预测缓存。}
在第$D$日，训练样本的历史预测起点依次取$\max(1,D-60),\ldots,D-1$，每个起点生成本日及次日目标，但仅保留目标日严格小于$D$的样本。这样，训练用到的实际结果均来自已经结束的完整日。岭回归保留全部144个日内点；提升树在拼接后的训练序列上每隔一个点取样，预测仍覆盖全部144个点。少于7天历史时，两种复杂模型均退回朴素基线。

基础模型每日拟合一次，并由同一午夜起点生成本日、次日及第三日共432步预测；第三日沿用同一模型和已知日历特征。这份缓存足以覆盖日内任一时刻发出的最长48小时查询。训练提前日数只有0和1，因此跨到第三日的部分是模型外推。每10分钟控制更新实际库存，使用当前已发布的光伏预报；基础负载和价格模型不在日内反复重训。

\paragraph{官方预报融合与历史误差块。}
每次只取允许发布时刻中最新且已经发布的一版光伏预报。其24个点对齐到发布时间后的第1至24小时，以最近已结束区间的实测光伏为插值锚点，形成10分钟预测；首个起点没有实测时使用附件1末点。超出该版预报覆盖范围后使用历史模型。消融实验中，未列入该实验的发布时刻对预测和购电调整同时不可用。

令当前绝对区间起点为$r=144D+s$，预测长度为$H$。历史块起点取$j=144d+s$，其中
\begin{equation}
 \max(0,D-60)\le d<D,\qquad j+H\le r.
\end{equation}
后一条件保证块内$H$个实际区间全部结束。历史基线取第$d$日原始午夜生成的完整三日预测，再截取$s$起的$H$步；若使用官方预报，则只叠加在历史起点$j$已经发布的版本。三条误差序列为负载和光伏的加性误差、价格的对数误差，按同一历史块配对，按正文情景公式转换为净负荷电量与正价格。

把合格块按时间排序，取等距索引$\operatorname{round}(\operatorname{linspace}(0,M-1,K))$，其中$M$为合格块数。12情境模型及问题三、四取$K=12$，等权使用选中的块；$M<K$时允许重复，无合格块时使用零残差。点预测消融$K=1$单独规定为零残差。若2月启动的模型尚无自己的1月预测记录，就在各历史起点严格重建朴素预测，并标记为因果朴素初始化；不会用2月拟合的模型回填1月预测。

\paragraph{模型选择与数值保护。}
12情境模型及问题三、四的三个候选各自从1月1日6000 kWh开始假想运行，以1月15日至31日的库存校正费用选模。它们1月15日的库存可以不同，比较时保存各自期初、期末值。具体规则是在费用不超过候选最小值的$1.005$倍的模型中，依朴素、岭回归、提升树的顺序取最简单者。其正式1月始终按朴素策略预热，2月继承该实际库存；候选实验不会改写这段运行。问题二风险裕度—即时反馈模型不使用这套选模结果。

负载和光伏预测截到非负区间。进入对数的价格设置$10^{-6}$下限；点预测和情景预测的对数价格在指数变换前截到$[-9,3]$，即价格约束于$[1.23\times10^{-4},20.09]$元/kWh。此界远宽于附件4的正常取值范围，用于防止外推造成数值溢出。任何非有限预测或回归系数均触发异常，不以替换结果的方式继续运行。

\subsection{参考末端、数值恢复与实际运行记录}
\label{app:records}

\paragraph{末端库存可达域。}
正文的短时域末端目标来自48小时参考轨迹。若实际库存偏离参考，原目标可能无法在剩余$H_M$步内达到，因此先将它投影到
\begin{equation}
 \left[\max\{1200,E_r-H_MB/\eta\},
       \min\{10800,E_r+H_MB\eta\}\right],\qquad B=5000/6\ \mathrm{kWh}.
\end{equation}
区间两端分别对应持续最大功率放电和充电，考虑了容量、功率与效率。投影后作为本次参考优化的硬末端条件。它只表达设备的库存可达性；实际安全层还受当期余电和缺口限制，可能无法跟踪参考末端，下一步应从实际$E$重新规划。2小时控制窗口更短，投影对避免功率限制导致的不可行尤为必要。本段只适用于12情境模型及问题三、四，不适用于问题二风险裕度—即时反馈模型。

\paragraph{求解结果数值复核。}
稀疏线性规划调用HiGHS，原始及对偶可行性容差均为$10^{-8}$；混合整数核验的相对间隙设为$10^{-8}$。程序要求求解状态为最优，并执行正文所述的同充同放消除，再由$C,D$重算整条$E$、各情景$X,W$及目标值。恢复后的参考约束最大偏差不得超过$2\times10^{-5}$，重算目标与求解器目标的差不得超过
$10^{-4}+10^{-7}\max\{1,|J_{\rm raw}|\}$。
这一步验证的是参考问题，实际电费仍由逐区间执行和结算产生。输入库存仅允许在边界外$10^{-7}$的舍入范围内截回容量界，更大越界直接报错。

\paragraph{运行记录内容。}
工程中，每个策略的\texttt{ledger.npz}保存全年逐区间数组，其字段含义如下。所有电量字段以kWh计，价格为元/kWh，费用为元。
\begin{center}
\small
\begin{tabularx}{\linewidth}{p{0.22\linewidth}Y}
\toprule
记录内容 & 字段与含义\\
\midrule
交易和实际供需 & \texttt{plan}为$G^0$，\texttt{effective}为$G$；\texttt{load,pv,price}为结算区间的实际负载电量、光伏电量和价格。\\
储能和补缺 & \texttt{charge,discharge,emergency,waste}对应$C,D,X,W$；\texttt{soc\_start,soc\_end}为两端$E$。\\
费用分项 & \texttt{base\_fee,up\_fee,down\_fee,emergency\_fee}记录原计划费、最终上调费、最终下调费、应急费；\texttt{cost}为其和。\\
控制和预测摘要 & 情景模型中\texttt{target\_charge,target\_discharge}为安全层执行前的首步建议；问题二风险裕度—即时反馈模型中二者记录即时反馈实际值。\texttt{expected\_net,expected\_price}为执行前参考预测；\texttt{issue\_index}为最近承诺更新时间。\\
\bottomrule
\end{tabularx}
\end{center}
\texttt{revisions.json.gz}记录每次调整时刻及调整后的剩余日承诺向量，保存到小数点后9位。\texttt{forecast\_diagnostics.json}保存日级训练审计、调用次数及冷启动记录数；\texttt{summary.json}保存策略配置和分评价窗口汇总。审计字段\texttt{training\_\allowbreak actual\_\allowbreak end\_\allowbreak exclusive}是可用历史的排他上界：索引达到该值的区间尚不能用于训练。\texttt{issue\_index}也不等于每次预测所用的完整预报元数据。

提交包将五种策略的完整账本保存为\texttt{q1\_ledger.npz}至\texttt{q4\_3\_ledger.npz}。各数组按10分钟顺序排列，首项对应2025年1月1日00:00--00:10；该格式保留全部字段并减少压缩包体积。提交包中的调整记录与预测审计也加相同策略前缀并以gzip压缩，例如\texttt{q2\_revisions.json.gz}与\texttt{q2\_forecast\_diagnostics.json.gz}。

\paragraph{在线检查与独立复核。}
预测返回时检查历史误差块结束位置不晚于当前起点、预报发布时间不晚于当前时刻。执行时即时检查库存边界、能量平衡及充放电与应急充电重叠；独立验证程序再从实际账本重算守恒、库存连续、功率界、安全层动作和各费用分项，物理容差为$10^{-6}$ kWh。工作簿另按原模板日期、时段及单元格逐项核对，连续应急区间先展开再与账本比较。

程序未保存全年每次预测的完整情景及全部信息元数据，也未持久化每次线性规划的完整目标记录。因此，信息边界结论由实现约束、即时断言、日级审计和测试共同支持；保存的账本足以重算实际物理量与电费，不能单独复原所有中间预测。真实组件的3日端到端测试将同一历史前缀后的负载、光伏、价格及尚未发布预报共同扰动，检查此前的计划、储能建议、实际库存和费用不变。指定日期的LP/MILP对照则是使用已实现净负荷的事后数值核验，不参与过去的运行决策。

\subsection{统计口径与信息下界}
\label{app:statistics}

\paragraph{预测评价指标口径。}
午夜24小时预测对正式期334天的全部日内点计算负载、光伏及价格MAE；问题二账本保存午夜联合预测，情景模型的每10分钟首步预测由账本中的情景均值计算净负荷电量及价格MAE、RMSE。附件3原始预报按发布时间与提前1至24小时分别统计，只使用实际光伏大于100 kW且存在实测的白天目标。后者剔除跨入2026年的无实测点。三者的提前量、样本与单位不同，图中的白天预报误差不能代替整日能量误差；光伏大量零值也使全时段MAPE不适合作为主要指标。

\paragraph{库存校正与配对区间估计。}
评价费用$C$之外，以固定$\lambda=\min p^{(\text{附件1})}/0.9$计算
$C^{\rm inv}=C+\lambda(E_{\rm start}-E_{\rm end})$。
两端库存同时计入，适用于正式期和1月候选比较；该校正不进入题目要求的实际电费。比较两种策略的日费用时，先按相同日期求差$\delta_d=C_d^{\text{前者}}-C_d^{\text{后者}}$，再做循环区块重采样：每次均匀抽取$\lceil334/7\rceil=48$个起点，各带连续7日并在序列末尾循环，拼接后截到334日。固定种子2026，重复2000次，取样本日均差的2.5\%和97.5\%分位数。正值表示后者节省，负值表示后者更贵。循环处理和单年季节性限制了区间的外推能力，该区间用于描述本年差异。

\paragraph{未来信息对照实验。}
仅价格全知实验用真实未来价格替换评价年内的预测价格，净负荷预测、参考结构及安全层保持原配置；超出2025年的参考部分仍使用预测价。它检验价格信息在这套近似策略中的作用，费用不必单调下降。

完全信息下界使用全年365天真实净负荷与相应真实价格，所有购电按基础价计费，初始$E=6000$，末端只受容量界约束。令其购电变量为$g_t$，保留相同储能物理约束并满足$g_t+v_t+D_t\ge\ell_t+C_t$。任一主结算实际策略都能取$g_t=G_t+X_t$形成可行解，且
\begin{equation}
 C-\sum_t p_t(G_t+X_t)
 =\sum_t p_t\left[\tfrac12(G_t-G_t^0)^+
                  +\tfrac12(G_t^0-G_t)^++4X_t\right]\ge0.
\end{equation}
因此该放松规划的最优值不超过相同全年范围的主策略费用。差额同时包含未来信息、交易限制和控制近似的影响；它不能与334日费用直接相减，也不能全部视为可实现的节省。

\subsection{校准、预测诊断与完整实验表}
\label{app:experiments}
下表保留候选校准和各变体的原始数值。33组实验包括问题二的无储能、无分位数裕度、70\%、90\%和12情境模型5组，问题三与问题四的点预测、储能、更新时点、情景数、控制时域、效率与结算解释等28组。参数变体分别运行并按自己的效率及结算条款核验；主参数和事前规则不因正式期结果而改变。
\begin{table}[H]\centering\small\caption{情景滚动模型一月后半月预测方法选择（储电量校正后的费用，元）}
\begin{tabular}{lrrrl}\toprule
策略 & 季节朴素 & 岭回归 & 梯度提升 & 入选方法\\\midrule
问题二12情境模型 & 993038.99 & 1021829.62 & 1007275.41 & 季节朴素\\
问题三 & 974841.35 & 994073.34 & 973691.02 & 季节朴素\\
问题四-2 & 1146794.67 & 1199761.83 & 1177412.12 & 季节朴素\\
问题四-3 & 1107397.57 & 1139949.50 & 1121541.43 & 季节朴素\\
\bottomrule\end{tabular}\end{table}
\begin{table}[H]\centering\small\caption{午夜对当天24小时的平均绝对预测误差}
\begin{tabular}{lrrr}\toprule
策略 & 负荷/kW & 光伏/kW & 电价/(元/kWh)\\\midrule
问题二 & 181.34 & 156.59 & 0.000000\\
问题三 & 282.38 & 197.49 & 0.000000\\
问题四-2 & 282.38 & 167.60 & 0.050916\\
问题四-3 & 282.38 & 197.49 & 0.050916\\
\bottomrule\end{tabular}\end{table}
{\small\begin{longtable}{lrrr}
\caption{33组对照与参数检验完整结果（二月至十二月）}\\\toprule
策略变体 & 实际费用/万元 & 紧急量/MWh & 库存校正费用/万元\\\midrule\endfirsthead
\caption[]{33组对照与参数检验完整结果（续）}\\\toprule
策略变体 & 实际费用/万元 & 紧急量/MWh & 库存校正费用/万元\\\midrule\endhead
问题二：无分位数裕度 & 1510.36 & 315.31 & 1510.37\\
问题二：无储能 & 1859.39 & 258.47 & 1859.39\\
问题二：70\%分位数 & 1456.17 & 136.19 & 1456.19\\
问题二：90\%分位数 & 1495.52 & 36.52 & 1495.55\\
问题二：12情境随机规划模型 & 1482.49 & 253.92 & 1482.52\\
问题三：原计划全额另加违约费 & 1484.12 & 165.99 & 1484.14\\
问题三：2小时 & 1468.53 & 178.81 & 1468.54\\
问题三：6小时 & 1465.22 & 180.46 & 1465.26\\
问题三：仅0时 & 1505.78 & 243.13 & 1505.81\\
问题三：0、12时 & 1474.24 & 197.08 & 1474.26\\
问题三：0、18时 & 1494.55 & 217.10 & 1494.57\\
问题三：0、6时 & 1489.54 & 198.05 & 1489.56\\
问题三：24情景 & 1462.12 & 171.67 & 1462.15\\
问题三：6情景 & 1483.99 & 223.53 & 1484.00\\
问题三：无储能 & 1897.66 & 257.66 & 1897.66\\
问题三：点预测 & 1624.09 & 833.48 & 1624.09\\
问题三：往返效率90\% & 1419.29 & 177.16 & 1419.33\\
问题四-2：价格全知 & 1540.94 & 267.45 & 1540.93\\
问题四-2：无储能 & 1973.43 & 319.20 & 1973.43\\
问题四-2：点预测 & 1869.85 & 1192.80 & 1869.82\\
问题四-3：原计划全额另加违约费 & 1557.73 & 163.61 & 1557.67\\
问题四-3：2小时 & 1539.54 & 176.11 & 1539.48\\
问题四-3：6小时 & 1538.07 & 177.72 & 1538.02\\
问题四-3：仅0时 & 1581.84 & 243.03 & 1581.78\\
问题四-3：0、12时 & 1547.26 & 194.42 & 1547.20\\
问题四-3：0、18时 & 1570.00 & 216.90 & 1569.95\\
问题四-3：0、6时 & 1563.37 & 196.79 & 1563.32\\
问题四-3：24情景 & 1529.55 & 169.15 & 1529.54\\
问题四-3：6情景 & 1562.38 & 219.05 & 1562.33\\
问题四-3：价格全知 & 1523.23 & 192.97 & 1523.22\\
问题四-3：无储能 & 1965.48 & 256.10 & 1965.48\\
问题四-3：点预测 & 1719.17 & 849.87 & 1719.15\\
问题四-3：往返效率90\% & 1491.10 & 172.31 & 1491.04\\
\bottomrule\end{longtable}}
逐时独立验证的最大能量平衡误差为$2.274\times 10^{-13}$ kWh；问题一LP与MILP费用差为0.000000000元。
全年365日完全信息松弛下界：固定电价13751266.38元，波动电价14494771.23元。该下界允许年末自由储电，与334日表内费用不能直接相减。
对应365日实际费用：问题二16757631.12元，问题三16684412.30元，波动电价下分别为18076656.90元、17650372.98元。差额同时包含信息、交易限制和策略近似的影响。
\resultfigure{forecast}{附件3光伏预报的平均绝对误差随提前小时数变化：正式334日，仅取实际光伏大于100 kW的时段}
\resultfigure{sensitivity}{情景数、控制范围及费用解释变化时，实际费用相对主策略的变化}
\FloatBarrier

